\problema{Word Break}{139}{src/07-dp/139-word-break.cpp}{M}

Dada una cadena \texttt{s} y un diccionario \texttt{wordDict}, determina si
\texttt{s} se puede partir en una secuencia de una o más palabras del
diccionario (las palabras se pueden reutilizar tantas veces como haga
falta).

\paragraph{La idea, con un ejemplo de la vida real.} Imagina que recibes un
mensaje de texto sin espacios, como \texttt{"leetcode"}, y tienes un
diccionario de palabras válidas. Quieres saber si es posible insertar
espacios en algunos lugares para que el mensaje quede formado completamente
por palabras del diccionario, por ejemplo \texttt{"leet" + "code"}. La
pregunta clave para resolverlo con \textbf{programación dinámica}
(\emph{dynamic programming}, DP) es: ¿el prefijo del mensaje hasta cierta
posición se puede partir bien? Si ya sabemos la respuesta para prefijos más
cortos, podemos construir la respuesta para prefijos más largos.

\begin{center}
\ttfamily
\begin{tabular}{l}
s = "leetcode", wordDict = ["leet", "code"]\\
"leetcode" = "leet" + "code" -> true\\
\end{tabular}
\end{center}

\paragraph{Cómo lo resuelve el código, paso a paso.}
\begin{enumerate}
  \item Guardamos el diccionario en un conjunto (\texttt{unordered\_set})
        para poder buscar una palabra en tiempo promedio constante.
  \item Creamos una tabla booleana \texttt{dp} de tamaño \texttt{n + 1},
        donde \texttt{dp[i]} indica si el prefijo de \texttt{s} de longitud
        \texttt{i} se puede partir completamente en palabras del
        diccionario. \texttt{dp[0] = true} (la cadena vacía siempre se puede
        ``partir'': caso base).
  \item Para cada posición \texttt{i} de \texttt{1} a \texttt{n}, probamos
        cada posible corte anterior \texttt{j} menor que \texttt{i}.
  \item Si \texttt{dp[j]} es verdadero (el prefijo hasta \texttt{j} ya se
        sabe partir bien) y el substring entre \texttt{j} e \texttt{i} está
        en el diccionario, entonces \texttt{dp[i]} también es verdadero:
        encontramos una forma válida de partir el prefijo de longitud
        \texttt{i}.
  \item La respuesta final es \texttt{dp[n]}, es decir, si la cadena
        completa se puede partir.
\end{enumerate}

\lstinputlisting{src/07-dp/139-word-break.cpp}

\complejidad{$O(n^2)$}{$O(n)$}

\paragraph{¿Qué significa esa complejidad?} Para cada una de las $n$
posiciones finales probamos hasta $n$ posibles puntos de corte, y comparar o
construir cada substring toma tiempo proporcional a su longitud, así que en
el peor caso el trabajo es cuadrático en el tamaño de la cadena. El espacio
extra es lineal, para la tabla \texttt{dp} (el diccionario en sí no cuenta
como espacio auxiliar del algoritmo, ya viene dado).

\paragraph{Consejos para la entrevista (Amazon).} La trampa más común es
asumir que el primer corte que funciona localmente (por ejemplo, tomar
siempre la palabra más larga posible desde el inicio) lleva a la solución
completa; con diccionarios que tienen palabras que se solapan, como
\texttt{\{"cats", "cat", "sand", "and", "dog"\}} sobre la cadena
\texttt{"catsandog"}, un corte greedy puede fallar aunque exista una
partición válida en otra combinación de palabras. Por eso se necesita
recordar, para CADA prefijo, si es alcanzable, en vez de intentar solo un
camino. También conviene mencionar la variante recursiva con memoización
como forma equivalente de llegar a la misma complejidad.
